\section{Part C: pointed torsors and relative moduli}
\label{sec:part-C}

\subsection{The distinction between scans 139 and 141}

We distinguish two assertions which the supplied study conflates.  Scan
29/139 states the injectivity assertion for admissible finite flat groups.
Scan 29/141 additionally requires the category to be closed under kernels
of pairs and under images, and explicitly says that all its groups are
admissible.  These additional hypotheses matter: the kernel of
$[x]\colon\alpha_p\to\alpha_p$ in the cusp example is not flat.  Thus that
example refutes scan 139, but does not refute scan 141 with its written
hypotheses.  None of the doubtful words in the supplied transcription
affects this distinction.

Here a finite flat group $G/S$ is \emph{admissible} if it has an exact
sequence of fppf sheaves
\[
1\longrightarrow G^0\longrightarrow G\longrightarrow\Gamma_G
\longrightarrow1
\]
with $G^0/S$ finite flat and radicial and $\Gamma_G/S$ finite \'etale.
For a finite flat group $G/S$, $Z_a(G)$ denotes isomorphism classes of
$G$-torsors over $S$ with a chosen lift of a fixed field-valued point
$a\colon\operatorname{Spec}\Omega\to S$.
All maps of pointed torsors preserve this point.  A class is minimal
relative to a category $\mathcal C$ if it is not induced from a proper
closed subgroup belonging to $\mathcal C$ through an allowed morphism.
The manuscript itself uses geometric marking: scan 29/126 explicitly
defines a marked point as an algebraically closed extension of a
residue field (the supplied K29-08 transcription, lines 80--84).
Some arguments below work for arbitrary field-valued marking, which
we indicate.  For the comparison with Nori we distinguish this
geometric convention from $\Omega=k$ with a $k$-rational lift on the
torsor.  A rational underlying point on $S$ does not by itself identify
the two functors.

\begin{theorem}[Proved: the literal assertion of 29/141]
\label{C:literal-141}
Let $S$ be connected, reduced and noetherian.  Let $G,G'/S$ be admissible
finite flat groups and let $u,v\colon G\rightrightarrows G'$ be group
homomorphisms whose equalizer $E$ is finite flat and admissible.  If a
pointed $G$-torsor $P$ has isomorphic pointed pushouts along $u$ and $v$,
then $P$ has a pointed reduction to $E$.

Consequently, if $\mathcal C$ is the category described on scan 29/141,
then, for every minimal $\alpha\in Z_a(G)$, the map
\[
\operatorname{Hom}_{\mathcal C}(G,G')\longrightarrow Z_a(G'),
\qquad u\longmapsto u_*\alpha,
\]
is injective.  Thus Proposition 4 is true with the literal hypotheses of
scan 141.
\end{theorem}

\begin{proof}
Let $G$ act on $G'$ by
$g\cdot z=u(g)zv(g)^{-1}$.  Its stabilizer at $1$ is $E$.  The finite
locally free quotient $G/E$ therefore has a closed orbit immersion
$G/E\hookrightarrow G'$.  Indeed the orbit map is a monomorphism of fppf
sheaves, and a finite monomorphism is a closed immersion.  We use here
the quotient theorem for a finite locally free equivalence relation,
applied to the right action of $E$ on the affine $S$-scheme $G$
\cite[Expos\'e~V, Theorem~4.1(iv)]{C_SGA3I}.

The maps $u,v$ induce maps of the finite \'etale component groups.
The component group of $E$ is the equalizer of these two maps: this may
be checked on algebraically closed geometric fibres, where the radicial
kernels have just one point.  Hence the support of $G/E$ is the inverse
image in $G'$ of the orbit of $1$ in the finite \'etale component set.
In particular, this support is an open and closed subscheme $U\subset G'$
and $G/E\hookrightarrow U$ is a nilpotent closed immersion.  The
nilpotence follows locally from noetherianity and equality of supports.

An isomorphism of the two pushout torsors is the same as a section
$s\colon S\to W:=P\mathbin{\times^G}G'$ of the associated finite scheme
for this action.  Its pointedness says that $s(a)$ belongs to the
associated orbit
\[
P\mathbin{\times^G}(G/E)=P/E\ \subset\ W.
\]
The associated open and closed subscheme
$P\mathbin{\times^G}U\subset W$ contains $s(a)$, so connectedness of
$S$ makes $s$ factor through it.  The remaining closed immersion
$P/E\hookrightarrow P\mathbin{\times^G}U$ is nilpotent.  Since $S$ is
reduced, $s$ factors through $P/E$.  Pulling back the $E$-torsor
$P\to P/E$ by this section gives the required pointed $E$-torsor.

If $P$ is minimal in $\mathcal C$ and $E\in\mathcal C$, this reduction
forces $E=G$.  The equalizer property then gives $u=v$.
\end{proof}

\begin{proposition}[Proved: an exact categorical criterion]
\label{C:criterion}
Fix a connected base and a category $\mathcal C$ of finite locally free
groups with the specified allowed group maps.  Suppose the category
contains the trivial group and zero homomorphisms, and has products
and equalizers represented by its scheme-theoretic group equalizers,
with their universal maps allowed.  For the pointed-torsor functor
$Z_a$, the following are
equivalent:
\begin{enumerate}
\item Proposition 4 holds for every minimal class in $Z_a$.
\item For every pair $u,v\colon G\rightrightarrows G'$ in $\mathcal C$
and $E=\operatorname{Eq}(u,v)$, the map
\[
Z_a(E)\longrightarrow
\{\alpha\in Z_a(G):u_*\alpha=v_*\alpha\}
\]
is surjective.
\end{enumerate}
If in addition every allowed homomorphism has its scheme-theoretic
image factorization in $\mathcal C$, with both the map onto the image
and the closed immersion allowed, these conditions are equivalent to
pro-representability of $Z_a$ by an inverse system of objects of
$\mathcal C$ with faithfully flat transition maps.
This is a necessary and sufficient criterion for an arbitrary such
triple of base, category, and allowed morphisms; the preceding theorem
verifies it under the actual manuscript hypotheses.
All these requirements on the factor maps hold for the $k$-defined
and all-$S$ homomorphism categories used below.  Membership of the
image object alone would not suffice for an arbitrary restricted
collection of allowed maps.
\end{proposition}

\begin{proof}
A proper closed subgroup of a finite locally free group over a connected
base has strictly smaller rank.  Therefore every pointed torsor has a
minimal reduction.  If (1) holds and $u_*\alpha=v_*\alpha$, choose a
minimal reduction $\beta\in Z_a(H)$ with inclusion $i\colon H\hookrightarrow
G$.  Then $(ui)_*\beta=(vi)_*\beta$, so $ui=vi$.  The equalizer property
factors $i$ through $E$, and pushing $\beta$ to $E$ proves (2).
Conversely, (2) reduces every equality for a minimal class to its
equalizer; minimality makes this equalizer all of $G$.

For pro-representability, use as indices the minimal pointed torsors.
There is at most one allowed group homomorphism between any two such
classes by (1).  A minimal reduction of the product of two pointed
torsors maps to both, so the index category is cofiltered.  A map between
minimal objects is faithfully flat: its scheme-theoretic image belongs
to $\mathcal C$, minimality of the target makes that image the target,
and its kernel is finite locally free by the equalizer hypothesis.
The quotient theorem then makes the map faithfully flat.  Every class
comes from a minimal reduction.  Two representing maps give the same
class precisely when they become equal after passage to a common
minimal refinement, by (1).  This proves the required colimit formula
for $Z_a$.

Conversely, suppose such a strict pro-object represents $Z_a$.  A minimal
class is represented at a finite stage by a homomorphism onto its target:
otherwise its scheme-theoretic image would give a proper reduction.
That homomorphism is faithfully flat by the kernel hypothesis.  Equality
of two pushed classes becomes equality at a later finite stage, and
faithful flatness then gives equality of the two target homomorphisms.
\end{proof}

\subsection{Normal bases and constant groups}

\begin{proposition}[Refuted: the broader assertion and the first candidate in L7]
\label{C:cusp-all-subgroups}
The assertion of scan 29/139 is false.  It remains false if minimality
is strengthened to prohibit reductions to every proper closed subgroup,
including nonflat ones, provided a torsor under such a subgroup means
an fppf-locally trivial torsor.
\end{proposition}

\begin{proof}
Let $k$ have characteristic $p>0$, let
$A=k[x,y]/(y^p-x^{p+1})=k[t^p,t^{p+1}]$, and mark the cusp of
$S=\operatorname{Spec}A$.  The finite flat $\alpha_p$-torsor
\[
P:\ z^p=x
\]
is nontrivial: in $K=k(t)$ its only possible section is $z=t$, whereas
$t\notin A$.  The two homomorphisms $0,[x]\colon\alpha_p\rightrightarrows
\alpha_p$ are distinct, but the pushout by $[x]$ has equation
$w^p=x^{p+1}=y^p$, so is trivial, with the required marked point.

We verify the stronger minimality statement.  A proper closed subgroup
$H\subset\alpha_{p,S}$ has trivial generic fibre: a subgroup of
$\alpha_{p,K}$ is either trivial or all of $\alpha_{p,K}$; in the
second case the defining ideal of $H$ vanishes generically and hence
vanishes because $\mathcal O(\alpha_{p,S})$ is torsion-free over $A$.
Suppose an fppf $H$-torsor reduced $P$.  On an fppf cover of $S$ that
trivializes this torsor, it gives local sections of $P$.  On every
double overlap the sections agree generically, because $H_K=1$.
Such an overlap is flat over the integral scheme $S$, so its generic
fibre is schematically dense: on affine pieces this is the injection
of a flat $A$-algebra into its localization by $A\setminus\{0\}$.
The sections therefore agree everywhere, since $P/S$ is separated.
They descend to a section of $P/S$, contradicting $t\notin A$.

Finally, the equalizer of $[x]$ and $0$ is
$\operatorname{Spec}A[w]/(w^p,xw)$.  It has generic rank $1$ and
closed-fibre length $p$, so is not flat.  This locates the exact
hypothesis missing from scan 139 and explains why the example does
not contradict Theorem~\ref{C:literal-141}.  The attached program checks
the polynomial identity and these lengths for $p=2,3,5,7$.
\end{proof}

\begin{proposition}[Refuted: C.1(i), even over a smooth affine surface]
\label{C:normal-counterexample}
Normality of $S$ does not make Proposition 4 true for all finite flat
$S$-groups and all $S$-homomorphisms, with minimality taken among finite
flat closed subgroup reductions.
\end{proposition}

\begin{proof}
Let $k$ be a field of characteristic $p>0$,
$A=k[x,y]$, $S=\operatorname{Spec}A$, and $a=(0,0)$.  This base is
regular, integral, connected, noetherian and normal.  Take
$G=\alpha_p^2$ and the pointed torsor
\[
P=\operatorname{Spec}
A[s,t]/(s^p-x^{p+1},\ t^p-xy^p),
\qquad (s,t)|_a=(0,0).
\]
The monic equations make $P/S$ finite locally free of rank $p^2$;
translation in $s,t$ makes it a $G$-torsor.  Define
$u=(y,-x)\colon G\to\alpha_p$ and $v=0$.  They are distinct, but
\[
(ys-xt)^p=y^px^{p+1}-x^pxy^p=0.
\]
Thus $u_*P$ and $v_*P$ are isomorphic as pointed torsors.

We verify minimality.  Put $K=k(x,y)$.  A proper nontrivial subgroup
of $\alpha_{p,K}^2$ has rank $p$ and is a linear copy of $\alpha_p$.
Indeed the equivalence with restricted Lie algebras identifies its
subgroups with linear subspaces of the two-dimensional restricted
Lie algebra having zero bracket and zero $p$-operation
\cite[Expos\'e~VIIA, Theorem~8.1.2]{C_SGA3I}.
If the generic torsor reduces to the line spanned by $(c,d)$, its class
is annihilated by $(d,-c)$.  Therefore
\[
(dx-cy)^p x\in K^p.
\]
Since $x\notin K^p$, this forces $dx=cy$.  Hence its only possible
rank-$p$ reduction has generic group the line spanned by $(x,y)$.
Its schematic closure in $G$ has coordinate algebra
\[
B=A[xz,yz]\ \subset\ A[z]/(z^p),
\qquad
B=\bigoplus_{j=0}^{p-1}(x,y)^jz^j.
\]
At the maximal ideal $\mathfrak m=(x,y)$ the vector-space dimension of
$B/\mathfrak m B$ is
\[
\sum_{j=0}^{p-1}(j+1)=\frac{p(p+1)}2>p,
\]
whereas its generic rank is $p$.  Thus $B$ is not flat at $\mathfrak m$.
Any finite flat subgroup with this generic fibre would equal its
schematic closure, because its coordinate algebra is torsion-free over
$A$.  No such subgroup exists.  The trivial subgroup cannot reduce
$P$, since $x^{p+1}\notin A^p$.  A closed subgroup of the same rank as
$G$ is $G$.  These exhaust the possible ranks, so $P$ is minimal.

The script \texttt{check\_counterexamples.py} verifies the pushout
identity and the fibre-dimension counts for $p=2,3,5,7$.  The argument
above proves them, and minimality, for every prime $p$.
\end{proof}

\begin{theorem}[Proved: C.1(ii)]
\label{C:k-morphisms}
Let $S$ be a connected reduced noetherian $k$-scheme, pointed by a
field-valued point.  Take finite $k$-group schemes as structure groups,
base changed to $S$, and only homomorphisms defined over $k$.  Proposition
4 holds, whether or not $S$ is proper or normal.
\end{theorem}

\begin{proof}
All finite groups over a field have the connected--\'etale exact
sequence, and its base change makes them admissible over $S$.
Equalizers of $k$-group homomorphisms are finite flat $k$-groups and
remain such after base change.  Theorem~\ref{C:literal-141} applies.

For an independent way to identify the disputed step on scan 142, form
the finite $k$-stack $W=[G'/G]$ for the action used in that theorem.
An isomorphism of pointed pushouts gives $S\to W$ whose marked fibre
is the object $1\in G'(k)$.  The scheme-theoretic image $\Gamma$ is
reduced and connected; its reducedness is checked after a smooth atlas
of $W$, since smooth base change preserves reducedness of $S$.
The marked fibre implies that the rational object $1$ factors through
$\Gamma$.  By Borne and Vistoli, Lemma 7.15(c), every morphism to a
reduced finite stack is flat.  Hence
$\operatorname{Spec}k\to\Gamma$ is a faithfully flat cover.  The
automorphism group of its object is exactly $E=\operatorname{Eq}(u,v)$,
since $\Gamma\hookrightarrow W$ is a closed immersion and therefore
fully faithful.  Thus $\Gamma\simeq BE$ and the factorization gives
the equalizer reduction.  This is a proof of the missing existence
step; it does not use the circular uniqueness assertion on scan 142
\cite[Lemma~7.15(c)]{C_BorneVistoli2015}.
\end{proof}

\begin{corollary}[Proved: the precise Nori comparison]
\label{C:nori}
If $k$ is perfect, $S/k$ is connected, reduced and of finite type, and
$a\in S(k)$, the pro-object for finite $k$-groups and
$k$-homomorphisms with a chosen \emph{$k$-rational} lift of $a$ on
each torsor is canonically Nori's fundamental group scheme.
With the manuscript's geometric-lift convention, the identification
is false in general, even when the underlying marked point is rational
and $S$ is proper and nonnormal.  It becomes true over algebraically
closed $k$, when these two conventions agree.
\end{corollary}

\begin{proof}
Theorem~\ref{C:k-morphisms} and Proposition~\ref{C:criterion} give its
universal bijection
\[
\operatorname{Hom}_k(\pi,G)\simeq Z_a(G)
\]
for every finite $k$-group $G$.  This is precisely the pointed universal
property in Nori, Chapter~II, \S1, Definition~1 and Propositions~1--2;
the original existence proposition already treats connected reduced
schemes without a properness assumption
\cite[Chapter~II, \S1, pp.~85--87]{C_Nori1982}.
More explicitly, reducedness over
the perfect field implies geometric reducedness.  Connectedness and
the rational point imply geometric connectedness: the finite set of
geometric connected components has a transitive Galois action and a
fixed component.  Borne--Vistoli, Proposition 5.5(b), Theorem 5.7 and
Remark 5.15 therefore construct Nori's group with exactly this
universal property.  Yoneda gives the canonical isomorphism of
profinite group schemes, not merely an identification on rational
points \cite[Proposition~5.5(b), Theorem~5.7, Remark~5.15]{C_BorneVistoli2015}.

For the distinction between conventions, first take
$S=\operatorname{Spec}\mathbb F_5$.  The rationally pointed functor
has only the trivial torsor under every finite group, so Nori's group
is trivial.  In contrast,
$\operatorname{Spec}\mathbb F_{25}\to\operatorname{Spec}\mathbb F_5$
is a nontrivial torsor under the constant group $\mathbf Z/2$, and
it has a chosen lift of the geometric point.

The same discrepancy occurs within the nonnormal scope of C.1(ii).
Let $S/\mathbb F_5$ be the projective cuspidal cubic $Y^2Z=X^3$,
marked at its rational smooth point at infinity.  Its normalization
$\mathbb P^1\to S$ is a finite universal homeomorphism: it is
bijective on geometric points, including the unique point over the
cusp.  Pullback therefore gives an equivalence on finite \'etale
covers \cite[Theorem~59.45.2, tag 04DZ]{C_StacksLimits}.
Over an algebraic closure, every connected \'etale cover of
$\mathbb P^1$ is trivial, by Riemann--Hurwitz.  Descent then identifies
finite \'etale covers of $\mathbb P^1_{\mathbb F_5}$, and hence of $S$,
with finite \'etale $\mathbb F_5$-schemes.  A rational lift at infinity
trivializes a torsor under such a group.  Thus the rationally pointed
functor for $\mathbf Z/2$ is a singleton, while the geometrically
pointed constant-field cover $S_{\mathbb F_{25}}\to S$ is nontrivial.
The script \texttt{check\_pointing\_and\_detection.py} checks the quadratic
finite-field algebra, its Frobenius involution and the absence of
$\mathbb F_5$-rational points on its spectrum.
\end{proof}

\begin{theorem}[Proved under the stated rational-point hypothesis]
\label{C:constant-all-maps}
Let $S$ be a connected reduced noetherian $k$-scheme with a point
$a\in S(k)$.  Let the objects be all base changes of finite $k$-groups,
but allow all $S$-group homomorphisms, including varying subgroup
embeddings.  Then the following is the complete answer for this
pointing convention:
\begin{enumerate}
\item If $\operatorname{char}k=0$, Proposition 4 always holds.
\item If $\operatorname{char}k=p>0$, Proposition 4 holds for all
minimal pointed torsors if and only if $\Gamma(S,\mathcal O_S)=k$.
\end{enumerate}
In particular nonnormality alone does not determine the answer.
Every proper geometrically connected reduced $k$-scheme with
$H^0(S,\mathcal O_S)=k$ satisfies the assertion, including singular
nonnormal examples.  The affine line in characteristic $p$ already
fails it.
\end{theorem}

\begin{proof}
Put $A=\Gamma(S,\mathcal O_S)$.  A morphism between two constant finite
group schemes has its coordinate coefficients in $A$.  Thus if $A=k$
all allowed morphisms descend to $k$, and
Theorem~\ref{C:k-morphisms} applies.  In characteristic zero all finite
$k$-groups are \'etale.  Their Hom functor is a finite \'etale
$k$-scheme.  A morphism from connected $S$ to it lands in a single
connected component; the rational point $a$ forces that component to
be $\operatorname{Spec}k$.  Again all allowed homomorphisms descend.

Now assume $\operatorname{char}k=p>0$ and $A\ne k$.  Let
$I=\ker(a^*\colon A\to k)$.  Then $I\ne0$ and $I\ne F(I)$.
Indeed, if Frobenius were surjective on $I$, every $h\in I$ would
have roots $h=h_n^{p^n}$ with $h_n\in I$.  At any point of $V(h)$,
all $h_n$ belong to the local maximal ideal.  Krull intersection in
the noetherian local ring gives $h=0$ in that stalk.  Hence $V(h)$
is open as well as closed.  It contains $a$, so connectedness forces
$h=0$ everywhere, contrary to $I\ne0$.

Choose $f\in I\setminus F(I)$ and $0\ne h\in I$.  The class of
$f$ defines a nontrivial pointed $\alpha_p$-torsor; its nontriviality
follows from the injection
$A/A^p\hookrightarrow H^1_{\mathrm{fppf}}(S,\alpha_p)$.
If there is a nonzero $r\in A$ with $r^pf\in A^p$, multiplication
by $r$ and the zero map have the same pushout of this torsor.  A
proper constant subgroup of $\alpha_p$ is trivial, so this is already
a minimal counterexample.

We may therefore suppose that
\begin{equation}
\label{C:no-scalar-annihilator}
r^pf\in A^p\quad\Longrightarrow\quad r=0.
\end{equation}
Take $G=\alpha_{p^2}\times\alpha_p$ and the pointed torsor
\[
P:\qquad s^{p^2}=h^pf,\qquad t^p=f,
\qquad (s,t)|_a=(0,0).
\]
The homomorphism $u=(F,-h)\colon G\to\alpha_p$ is nonzero and
kills this class, because $(s^p-ht)^p=0$.  We show that $P$ is minimal
even when subgroup embeddings are allowed to vary over $S$.

Suppose it reduces to a proper constant subgroup $H_0\times_k S$.
Its rank is at most $p^2$.  Evaluating its embedding at $a$ identifies
$H_0$ with a subgroup of $\alpha_{p^2,k}\times\alpha_{p,k}$.
If $F_{H_0}=0$, its projection into the first factor lies in
$\alpha_p$, so the Frobenius pushout of the first component of $P$
is trivial.  This contradicts
$h^pf\notin A^p$, by~\eqref{C:no-scalar-annihilator}.
Otherwise $H_0$ has rank $p^2$ and its first projection at $a$ is
an isomorphism onto $\alpha_{p^2}$: a proper subgroup of
$\alpha_{p^2}$ lies in $\alpha_p$, which would make both coordinate
projections have zero Frobenius.  Hence $H_0\simeq\alpha_{p^2}$.

Over a reduced base, every homomorphism
$\alpha_{p^2}\to\alpha_{p^2}\times\alpha_p$ has the form
\[
z\longmapsto (b_0z+b_1z^p,\ c z^p),
\qquad b_0,b_1,c\in A.
\]
For a closed immersion $b_0$ must be a unit: on any fibre where it
vanishes, the whole $\alpha_p$ kernel of Frobenius is killed.
The first coordinate is consequently an automorphism, so the image
is the graph of $dF$ for some $d\in A$.  A reduction to this graph
would imply
\[
f-d^ph^pf=(1-dh)^pf\in A^p.
\]
But $(1-dh)(a)=1$, so $1-dh\ne0$, contradicting
\eqref{C:no-scalar-annihilator}.  This proves minimality and the
failure of Proposition 4.

For an explicit instance take $S=\mathbb A^1_k$, $a=0$ and $f=h=x$.
Then $r^px\in k[x]^p$ implies $r=0$, by the exponents modulo $p$.
The kernel of $u$ is the finite flat group
\[
H_x:\quad a^p=xb,\qquad b^p=0.
\]
It is $\alpha_{p^2}$ where $x$ is invertible and
$\alpha_p^2$ at $x=0$.  It is not a constant group scheme, explaining
exactly why the category in this example fails the kernel-closure
condition of scan 141.
\end{proof}

\begin{corollary}[Proved without a rational-point assumption]
\label{C:nonfield-all-maps}
Let $k$ have positive characteristic, and let $S$ be connected, reduced
and noetherian, with any geometric marking.  If
$A=\Gamma(S,\mathcal O_S)$ is not a field, Proposition 4 fails for
constant finite $k$-groups with all $S$-homomorphisms.  If $A=k$, it
holds.  Thus the only case not covered by these two tests has global
functions a field properly extending $k$; it is covered by the exact
criterion in Theorem~\ref{C:all-constant-classification} below.
\end{corollary}

\begin{proof}
Choose a nonzero nonunit $h\in A$ and a point $b\in S$ where it
vanishes; such a point exists, since a nowhere-vanishing global
function is a unit.  Put $I=\ker(A\to\kappa(b))$.  The Krull
intersection argument in the preceding proof gives
$I\setminus F(I)\ne\varnothing$.  Choose $f$ there and use exactly
the scalar-annihilator dichotomy and the torsor
$s^{p^2}=h^pf$, $t^p=f$ from that proof.

Only one descent detail changes.  A hypothetical subgroup $H_0$
having nonzero Frobenius is now identified with $\alpha_{p^2}$ after
extension to $\kappa(b)$.  Every field form of $\alpha_{p^2}$ is
trivial.  To see this, its automorphism group consists of
$z\mapsto b_0z+b_1z^p$ with $b_0$ invertible and has a split exact
sequence with kernel $\mathbb G_a$ and quotient $\mathbb G_m$.
The vanishing of $H^1(k,\mathbb G_a)$ and $H^1(k,\mathbb G_m)$
therefore gives the triviality of its fpqc forms.  Hence
$H_0\simeq\alpha_{p^2,k}$, and the graph argument applies unchanged:
$(1-dh)(b)=1$ yields the contradiction.  The torsors and their
pushouts have unique points over any algebraically closed field, so
the construction respects the originally prescribed geometric
marking, even when it differs from $b$.

If $A=k$, all morphisms between constant finite group schemes have
coefficients in $k$, so Theorem~\ref{C:k-morphisms} applies.
\end{proof}

\begin{proposition}[Refuted: a field of global functions is not sufficient]
\label{C:field-descent}
Without rational pointing over $k$, Proposition~4 for constant finite
$k$-groups and all $S$-homomorphisms can fail even when $S$ is proper,
geometrically integral over its field of global functions, and
non-normal.  There are examples in characteristic zero and over
finite fields.  In particular, the rational-point qualification in
Theorem~\ref{C:constant-all-maps} cannot be omitted.
\end{proposition}

\begin{proof}
First take
\[
k=\mathbb Q,\qquad L=\mathbb Q(\sqrt2),\qquad
b=2+\sqrt2,\qquad M=L(\sqrt b).
\]
The extension $M/k$ is cyclic of degree four.  Explicitly, if
$r=\sqrt b$, then $r^4-4r^2+2=0$, an Eisenstein polynomial at $2$,
and $\sigma(r)=r^3-3r$ defines an automorphism with
$\sigma^2(r)=-r$ and $\sigma(\sqrt2)=-\sqrt2$.
Let $G_0/k$ be the finite \'etale group whose geometric points are
$\mathbb F_3^2$ and on which this Galois generator acts by
\[
T=\begin{pmatrix}0&-1\\1&0\end{pmatrix}.
\]
This is completely explicit finite Galois descent: its coordinate
algebra is the algebra of functions $c\colon\mathbb F_3^2\to M$
satisfying $c(Tv)=\sigma(c(v))$, with comultiplication induced by
addition of vectors.  As an algebra it is $k\times M\times M$.

Over $L$, the action is $T^2=-1$.  Every order-three subgroup of
$G_{0,L}$ therefore has the quadratic character of $M/L$; denote
this $L$-group by $H$.  No such group descends to $k$.  Indeed a
$k$-form of the cyclic group of order three is given by a quadratic
character, with associated extension $N/k$ of degree at most two.
If its restriction were the character of $M/L$, we would have
$NL=M$.  This would make $M/k$ biquadratic or of degree at most two,
contrary to its cyclic degree-four Galois group.  Thus $G_{0,L}$
has no nontrivial proper subgroup isomorphic to the base change of
a finite $k$-group, even though $G_{0,L}\simeq H\times H$.

The elliptic curves
$E_b\colon y^2=x^3+b$ and $E_{-27b}$ are smooth.  The separable
isogeny of degree three
\[
\phi\colon E_b\longrightarrow E_{-27b},\qquad
(x,y)\longmapsto
\left(\frac{x^3+4b}{x^2},\frac{y(x^3-8b)}{x^3}\right)
\]
has kernel $\{O,(0,\sqrt b),(0,-\sqrt b)\}=H$.  These are the
V\'elu formulas for this subgroup; they are defined over $L$ because
the subgroup is Galois invariant
\cite[\S3, equations (8)--(11), Remark~3]{C_Velu1971}.
The point $(0,\sqrt b)$ has order three since its tangent line
meets the cubic with multiplicity three.  The displayed formula also
directly verifies the target equation, extends over the smooth
projective curves, and has degree three.  Its source is geometrically
connected, so it is a nontrivial pointed $H$-torsor.

Let $D/L$ be the projective cuspidal cubic $Y^2Z=X^3$, choose its
smooth point $d=[0:1:0]$, and put
\[
S=E_{-27b}\times_L D,
\qquad a=(O,d).
\]
We regard $S$ as a $k$-scheme and use this $L$-valued point, or a
geometric point above it, as marking.  The scheme is noetherian,
proper, connected and reduced, with
$\Gamma(S,\mathcal O_S)=L$.  It is non-normal: base changing the
finite normalization of the cusp by the smooth curve gives a finite
birational map which is not an isomorphism, as is already seen on
the fibre above $O$.  Pull back $\phi$ along the first projection
to obtain $Q/S$.  It remains nontrivial because restriction to
$E_{-27b}\times\{d\}$ is $\phi$.

Embed $H$ as the first factor of $G_{0,L}=H\times H$ and set
$P=Q\times^H G_{0,S}$.  Projection to the first factor recovers
$Q$, so $P$ is nontrivial.  Any proper constant subgroup object
$H_0\times_k S\subset G_{0,S}$ would, after evaluation at $a$,
give a subgroup $H_{0,L}\subset G_{0,L}$.  Its order is $1$ or
$3$; order $3$ was ruled out above.  Consequently $P$ is minimal
among the objects allowed in C.1(iii).  The nonzero projection
$u\colon H\times H\to H\times H$ onto the second factor kills
$P$, while the zero homomorphism also gives the trivial pointed
torsor.  Thus the required injectivity fails.

For a positive-characteristic example, use exactly the same
construction with $k=\mathbb F_5$, $L=\mathbb F_{25}$,
$M=\mathbb F_{625}$ and $b=u$, where $u^2=2$ in $L$.
Here $u^{12}=-1$, so $u$ is not a square and $M=L(\sqrt u)$.
The Galois group of $M/k$ is again cyclic of order four.
All the preceding subgroup, isogeny and torsor arguments apply.

The attached script \texttt{check\_field\_descent.py} checks all nine
vectors and all four subgroup lines, the cyclic-quartic algebra
identities, all twenty-five elements of $\mathbb F_{25}$, and the
cleared isogeny equation.  This also verifies concretely that the
obstruction concerns descent of subgroup objects, not non-normality
of $S$ or variation of global functions along $S$.
\end{proof}

\subsection{The exact descent criterion for every base in C.1(iii)}

\begin{lemma}[Proved: finite targets detect every proper subgroup]
\label{C:detect-proper-subgroup}
Let $L/k$ be an arbitrary extension of fields, let $G_0$ be a finite
$k$-group scheme, and let $J\subsetneq G_{0,L}$ be a closed subgroup.
There exist a finite $k$-group scheme $Q_0$ and distinct
$L$-homomorphisms
\[
u,v\colon G_{0,L}\rightrightarrows Q_{0,L}
\quad\hbox{such that}\quad u|_J=v|_J.
\]
Neither perfectness nor separability of $L/k$ is required.
\end{lemma}

\begin{proof}
Let $E_0=\pi_0(G_0)$ be the maximal finite \'etale quotient, and let
$J_e\subset E_{0,L}$ be the image of $J$.  The connected--\'etale
sequence exists over any field; after passage to an algebraic closure
its reduced subgroup gives a section of the component quotient
\cite[Chapter~IV, Proposition~4.45]{C_MoonenQuotients}.

Suppose first that $J_e\ne E_{0,L}$.  The scheme
$W_0=E_0\amalg E_0$ is finite \'etale over $k$, and its scheme of
automorphisms $Q_0=\operatorname{Aut}_k(W_0)$ is a finite \'etale
$k$-group.  Indeed after a finite separable splitting extension this
is the constant symmetric group on $2\operatorname{rank}(E_0)$
letters, with its natural descent datum.  Let
$r\colon E_0\to Q_0$ act by left translation on both copies.
Define $\tau\in Q_0(L)$ to exchange the two copies over the open
and closed subset $J_e\subset E_{0,L}$ and to fix their complements.
Left translation by $J_e$ preserves $J_e$ and its complement, so
$r_L|_{J_e}$ commutes with $\tau$.  Left translation by a geometric
point outside $J_e$ does not commute with $\tau$: applying the two
compositions to the identity in the first copy gives points in
different copies.  Thus $r_L$ and $\tau r_L\tau^{-1}$ agree on
$J_e$ and are distinct.  Composing them with
$G_{0,L}\to E_{0,L}$ proves the assertion in this case.

Suppose now that $J_e=E_{0,L}$.  Since $J$ is proper this case has
positive characteristic $p$.  Let $V_0$ be the left regular
representation of $G_0$ on $k[G_0]$; explicitly
$(g\cdot f)(x)=f(g^{-1}x)$.  The quotient $J\backslash G_{0,L}$
is a finite $L$-scheme of length
$\operatorname{rank}(G_0)/\operatorname{rank}(J)>1$, and
\[
L[J\backslash G_{0,L}]=V_L^J,\qquad V_L^{G_{0,L}}=L.
\]
The quotient and this invariant-algebra description follow from the
finite locally free quotient theorem
\cite[Expos\'e~V, Theorem~4.1(iv)]{C_SGA3I}.
We may therefore choose $w\in V_L^J\setminus V_L^{G_{0,L}}$.
The map of $L$-schemes
\[
c\colon G_{0,L}\longrightarrow V_L,
\qquad c(g)=g\cdot w-w,
\]
is a nonzero cocycle, with
$c(gh)=c(g)+g\cdot c(h)$, and vanishes on $J$.

Over an algebraic closure $\overline L$, every connected component
of the finite scheme $G_{0,\overline L}$ has one reduced point.
The equality $J_e=E_{0,L}$ says that $J_{\overline L}$ meets all
these components.  It therefore contains
$(G_{0,\overline L})_{\mathrm{red}}$ as a closed subscheme.
Consequently every coordinate function of $c$ vanishes on the
geometric reduction of $G_{0,L}$.  These functions are nilpotent in
the finite $L$-algebra $L[G_0]$.  We may choose an integer $m$ such
that their $p^m$-th powers all vanish.

Let
\[
N_0=\ker\bigl(F^m_{V_0/k}\colon V_0\to V_0^{(p^m)}\bigr).
\]
This is a finite $k$-group, isomorphic as a group scheme to
$\alpha_{p^m}^{\dim_k V_0}$.  It is preserved by the $G_0$-action:
relative Frobenius is natural for the linear action, so it carries
the action on a point killed by Frobenius to zero.  The preceding
coordinate calculation says exactly that $c$ factors through
$N_{0,L}$.  We may form the finite $k$-group
$Q_0=N_0\rtimes G_0$.  The two maps
\[
g\longmapsto(c(g),g),\qquad g\longmapsto(0,g)
\]
from $G_{0,L}$ to $Q_{0,L}$ are homomorphisms by the cocycle
identity, agree on $J$, and are distinct.  This proves the lemma.
\end{proof}

We now give a criterion in terms of a specified profinite group and
descent of its finite quotient images.  It does not quantify over
pairs of pushout homomorphisms, and the preceding lemma proves that
there is no further obstruction hidden by the available target
groups.

\begin{theorem}[Proved: a complete intrinsic criterion for C.1(iii)]
\label{C:all-constant-classification}
Let $S$ be a connected reduced noetherian $k$-scheme, with its fixed
geometric point $a$, and put $A=\Gamma(S,\mathcal O_S)$.  The objects
of $\mathcal C$ are all finite groups $G_0\times_k S$, and its maps
are all $S$-group homomorphisms.  The following alternatives give a
necessary and sufficient criterion for Proposition~4 for every such
$S$, including every nonnormal $S$.
\begin{enumerate}
\item If $\operatorname{char}k=p>0$ and $A$ is not a field, the
assertion is false.
\item If $\operatorname{char}k=p>0$ and $A=L$ is a field, use this
field $L$ in the criterion below.
\item If $\operatorname{char}k=0$, let $L\subset A$ be the field
of elements algebraic over $k$, and use this $L$ in the criterion
below.  There is no requirement that $A$ itself be a field.
\end{enumerate}
In alternatives (2) and (3), view $S$ as an $L$-scheme and let
$\Pi=\pi^{\mathrm{fin},L}_{S,a}$ be the pro-finite $L$-group
representing geometrically pointed torsors under \emph{all} finite
$L$-groups and $L$-homomorphisms.  Its existence follows from
Theorem~\ref{C:k-morphisms} and Proposition~\ref{C:criterion}.
Then Proposition~4 holds if and only if the following condition holds:

\smallskip
\noindent\emph{For every finite $k$-group $G_0$ and every proper
closed $L$-subgroup $J\subsetneq G_{0,L}$ that is a faithfully flat
finite quotient of $\Pi$, there exists a proper closed subgroup
$H\subsetneq G_{0,L}$ containing $J$ such that
$H\simeq H_0\times_k L$ for a finite $k$-group $H_0$.}
\smallskip

The isomorphism in this condition is an isomorphism of group schemes;
its embedding into $G_{0,L}$ may be defined over $L$.  In particular,
the condition is automatic if $L=k$.  Equivalently, the condition says
that every torsor minimal among the constant $k$-group objects has
full monodromy in the category of all finite $L$-groups.  It is an
exact descent criterion, not a claim that the field $L$ alone
determines the answer.
\end{theorem}

\begin{proof}
Alternative (1) is Corollary~\ref{C:nonfield-all-maps}.  We justify
all the comparisons needed for the other alternatives.

In alternative (2), a map of the constant finite group schemes has
its coordinate coefficients in $A=L$.  Hence for every pair of
finite $k$-groups
\[
\operatorname{Hom}_S(G_{0,S},G'_{0,S})
=\operatorname{Hom}_L(G_{0,L},G'_{0,L}).
\]
In characteristic zero, every finite $k$-group is \'etale and the
scheme of its homomorphisms to another finite group is finite
\'etale.  Each map $S\to\operatorname{Hom}_k(G_0,G'_0)$ therefore
has all its coefficients algebraic over $k$.
To check that these coefficients lie in a field, note that $A$ is
reduced and has no nontrivial idempotents.  Any subalgebra generated
by finitely many elements algebraic over $k$ is finite and reduced,
and hence is a product of finite separable field extensions of $k$.
The absence of nontrivial idempotents forces it to be one field.
Taking their directed union proves that $L$ is a field and gives the
same displayed Hom equality in alternative (3).

The structural map $S\to\operatorname{Spec}L$ is faithfully flat
and quasi-compact: flatness holds over a field, nonemptiness gives
surjectivity, and $S$ is noetherian.  Thus an $L$-homomorphism of
finite groups is a closed immersion if and only if its pullback to
$S$ is one.  Consequently the subgroup objects allowed in
$\mathcal C$ are exactly the subgroups $H\subset G_{0,L}$ whose
group scheme has a $k$-form, with arbitrary $L$-defined embeddings.
This comparison is necessary; replacing them by all $L$-subgroups
would change minimality.

For a pointed torsor class $\alpha\in Z_a(G_{0,S})$, the universal
property of $\Pi$ gives
$\rho_\alpha\colon\Pi\to G_{0,L}$.  Its schematic image is a finite
closed subgroup $J$, and $\Pi\to J$ is faithfully flat.  Indeed the
map factors through a finite stage of the strict pro-system defining
$\Pi$, and the image of a homomorphism of finite groups over a
field is its faithfully flat quotient by the kernel.  Naturality of
the universal property gives both of the following equivalences:
\[
\begin{aligned}
\alpha\text{ reduces to }H
&\ \Longleftrightarrow\ \rho_\alpha\text{ factors through }H
\ \Longleftrightarrow\ J\subset H,\\
u_*\alpha=v_*\alpha
&\ \Longleftrightarrow\ u\rho_\alpha=v\rho_\alpha.
\end{aligned}
\]

Assume the displayed descent condition.  If $\alpha$ is minimal
among the allowed constant objects and $J$ were proper, that
condition would give a proper allowed subgroup $H$ containing $J$,
contradicting the first equivalence.  Thus $J=G_{0,L}$, so
$\rho_\alpha$ is faithfully flat.  The second equivalence then
implies $u=v$.  This proves Proposition~4.

Conversely, if the descent condition fails, choose a proper subgroup
$J\subset G_{0,L}$ and a faithfully flat map $\Pi\to J$ witnessing
its failure.  Their composite defines a pointed $G_{0,S}$-torsor
class $\alpha$.  The first equivalence and the absence of a proper
allowed overgroup of $J$ show that $\alpha$ is minimal in
$\mathcal C$.  Lemma~\ref{C:detect-proper-subgroup} gives distinct
maps $u,v\colon G_{0,L}\rightrightarrows Q_{0,L}$ with a finite
$k$-defined target and with $u|_J=v|_J$.  They are allowed
$S$-homomorphisms, and the second equivalence gives
$u_*\alpha=v_*\alpha$.  This is exactly the required counterexample
to Proposition~4.  Both directions of the criterion are proved.
\end{proof}

\subsection{Relative representability: the \'etale case}

\begin{theorem}[Proved under finite-presentation hypotheses]
\label{C:etale-moduli}
Let $S$ be locally noetherian and let $f\colon X\to S$ be proper,
flat and of finite presentation, with a section $\sigma$ and
$f_*\mathcal O_X=\mathcal O_S$ universally.  Let $G/S$ be finite
\'etale.  The pointed-torsor functor is represented by a quasi-finite
\'etale algebraic space over $S$.  It need not be separated.
If, in addition, the fibres of $f$ are geometrically reduced, the
representing space is a separated quasi-finite \'etale scheme.
\end{theorem}

\begin{proof}
Consider the stack
\[
\mathcal M_G=
\underline{\operatorname{Hom}}_S(X,BG)
\mathbin{\times}_{BG,\,\sigma^*}S,
\]
where $S\to BG$ is the trivial torsor.  Its objects over $T$ are
$G_T$-torsors on $X_T$ with a trivialization along $\sigma_T$.
Hall--Rydh, Theorem 1.2, gives algebraicity, local finite presentation
and affine diagonal: the target $BG$ has finite, hence affine,
diagonal \cite[Theorem~1.2]{C_HallRydh2019}.

Every pointed automorphism is trivial.  Indeed its equalizer with
the identity is open and closed in $X_T$, since the torsor is finite
\'etale; every geometric fibre of $X_T/T$ is connected and contains
the section, so this equalizer is all of $X_T$.  Thus $\mathcal M_G$
is an algebraic space and represents the class functor itself.
Finite \'etale covers, their group actions, and their trivializations
are invariant under nilpotent thickenings.  Hence $\mathcal M_G/S$
is formally \'etale, and therefore \'etale.

For quasi-finiteness we use the constructibility statement for proper
direct images of constructible stacks: SGA 1, XIII, Corollary 6.2,
applied to the stack of $G$-torsors.  Its rigidification along the
section is a constructible sheaf of finite sets.  The corresponding
\'etale algebraic space is of finite type and has finite geometric
fibres, hence is quasi-finite
\cite[Expos\'e~XIII, Corollary~6.2]{C_SGA1}.

Assume now that the fibres are geometrically reduced.  Let $T$ be a
trait over $S$, with fraction field $K$, and let $P,Q$ be two pointed
torsors over $X_T$ with a pointed generic isomorphism.  The scheme
$I=\operatorname{Isom}_{G}(P,Q)$ is finite \'etale over $X_T$.
Consequently $I/T$ is proper, flat and has geometrically reduced
fibres.  Its Stein scheme $D=\operatorname{Spec}(g_*\mathcal O_I)$
is finite \'etale over $T$, by SGA~1, X, Proposition~1.2; its
geometric fibres index the connected components of the fibres of
$I/T$.  The given generic isomorphism is a section $X_K\to I_K$.
Its image is an open and closed component isomorphic to the
geometrically connected $X_K$, so it specifies a $K$-point of $D$.
Because $D/T$ is finite and $T$ is normal, this $K$-point extends
uniquely to a section of $D/T$.  Pulling back $I\to D$ gives an
open and closed subscheme $J\subset I$ whose generic fibre is that
chosen component.  The finite \'etale map $J\to X_T$ has degree
$1$ on the generic fibre.  Its degree is locally constant on the
connected $X_T$, so it is an isomorphism.  This extends the generic
isomorphism $P_K\simeq Q_K$.  Along $\sigma_T$ it preserves the
pointings because two sections of a finite \'etale separated
$T$-scheme that agree generically agree everywhere.

This proves the valuative separatedness criterion; for the
locally noetherian, quasi-finite, locally finite type space in
question traits suffice.  A separated quasi-finite algebraic space
over a scheme is a scheme
\cite[Expos\'e~X, Proposition~1.2]{C_SGA1}.
\end{proof}

\begin{proposition}[Refuted: separatedness in C.2 as stated]
\label{C:etale-nonseparated}
Even with $S$ a regular trait and $G=\mathbb Z/2\mathbb Z$ constant,
the hypotheses in C.2 do not force a separated representing object.
\end{proposition}

\begin{proof}
Let $k$ be algebraically closed of characteristic different from $2$,
and choose a connected \'etale double cover $q\colon Q_0\to C$ of a
smooth proper connected curve.  Write $\tau$ for its deck involution.
Set
\[
A=k[t]_{(t)},\quad
E=\operatorname{Spec}A[z]/(z^2-t^2),\quad
X=(Q_0\times_k E)/\langle(\tau,z\mapsto-z)\rangle.
\]
The diagonal involution is free because $\tau$ has no geometric fixed
point.  Its quotient exists as a projective scheme, and
$P=Q_0\times E\to X$ is a finite \'etale $\mathbb Z/2$-torsor.
The morphism $X\to\operatorname{Spec}A$ is proper and flat: these
properties descend through the finite \'etale surjective map $P\to X$,
and $P/A$ is proper and flat.

For every $A$-algebra $R$,
\[
H^0(P_R,\mathcal O)=R[z]/(z^2-t^2),\qquad
H^0(X_R,\mathcal O)
=\bigl(R[z]/(z^2-t^2)\bigr)^{z\mapsto-z}=R.
\]
The first equality uses $H^0(Q_0,\mathcal O)=k$ and base change over
a field.  The second uses $2\in R^\times$ and the basis $1,z$.
These calculations sheafify, proving universal cohomological
flatness in dimension zero.

Choose $q_0\in Q_0(k)$.  The map $t\mapsto(q_0,z=t)$ into $P$ gives
a section $\sigma$ of $X/A$ and a pointing of $P$.  Over
$K=k(t)$ the two factors $z-t$ and $z+t$ are coprime.  The two copies
of $Q_{0,K}$ in $P_K$ are interchanged by the involution, so
$X_K\simeq Q_{0,K}$ and $P_K$ is the trivial pointed double torsor.
On the special fibre, however, reduction gives the nontrivial
connected cover $Q_0\to C$.  Hence $P$ is not the trivial torsor.

If the functor were represented by a separated algebraic space, the
two sections defined by $P$ and the trivial pointed torsor would agree
on the trait whenever they agreed at its generic point.  They do not.
This contradicts separatedness.  All the hypotheses of C.2, including
the section and universal equality of direct images, have been
verified; the special fibre is deliberately nonreduced.
\end{proof}

\subsection{Finite flat groups: the class functor need not be a sheaf}

\begin{theorem}[Proved: the general representing stack]
\label{C:finite-flat-stack}
Under the finite-presentation hypotheses of
Theorem~\ref{C:etale-moduli}, let $G/S$ be finite locally free, without
an \'etaleness assumption.  The groupoid of pointed $G$-torsors is
represented by the algebraic stack
\[
\mathcal M_G=
\underline{\operatorname{Hom}}_S(X,BG)
\mathbin{\times}_{BG}S,
\]
locally of finite presentation with affine diagonal.  Its functor of
isomorphism classes need not be an fppf sheaf and in general has no
representing scheme or algebraic space.

If $G$ is commutative, the stack has trivial inertia, is an algebraic
space, and its class functor is naturally
\[
\underline{\operatorname{Hom}}_{S\text{-groups}}
(G^D,\operatorname{Pic}_{X/S}).
\]
Neither \'etaleness nor quasi-finiteness follows in this commutative
case.
\end{theorem}

\begin{proof}
The stack assertion is Hall--Rydh, Theorem 1.2, followed by the
displayed fibre product.  For commutative $G$, an automorphism of a
$G$-torsor is a section of $G_{X_T}$.  Since $G_T$ is affine and
$f_{T*}\mathcal O_{X_T}=\mathcal O_T$, every such section comes from
$G(T)$.  Pointing makes it the identity.  Thus the stack is an
algebraic space.

Raynaud's correspondence identifies the sheaf
$R^1 f_{\mathrm{fppf}*}G$ with
$\underline{\operatorname{Hom}}(G^D,\operatorname{Pic}_{X/S})$;
we use the original Proposition~6.2.1(1), with $M=G^D$ and $M'=G$.
The stated universal direct-image hypothesis gives its required
equality on the fppf site
\cite[Proposition~6.2.1(1), pp.~50--51]{C_Raynaud1970}.
The explicit pointed form also appears in Mehidi--Poiret, Theorem~2.2
and Remark~2.3 \cite{C_MehidiPoiret2025}.
The section splits pullback from the base, and a pointed torsor is
its uniquely rigidified relative class.  This proves the formula.

For example, if $X$ is a supersingular elliptic curve over an
algebraically closed field of characteristic $p$ and $G=\alpha_p$,
the Hom group is $\operatorname{Hom}(\alpha_p,X)$, represented by
$\mathbb G_a$.  This holds on arbitrary, including nonreduced, test
schemes $T$: naturality of relative Frobenius makes every homomorphism
$\alpha_{p,T}\to X_T$ factor through
$\ker(F_{X_T/T})\simeq\alpha_{p,T}$.  A primitive element of
$\mathcal O_T[z]/(z^p)$ has the form $bz$, so
$\underline{\operatorname{End}}(\alpha_p)=\mathbb G_a$ on all such
tests.  This is not
quasi-finite or \'etale.  Taking $G=\mu_p$ instead gives $X[p]$,
which is finite but non-\'etale.

The claimed failure of the sheaf condition for a noncommutative finite
group is proved by the next proposition.
\end{proof}

\begin{proposition}[Refuted: representability of the general class functor]
\label{C:non-sheaf}
There are a proper pointed scheme $X/k$ with
$H^0(X,\mathcal O_X)=k$ universally and a finite flat noncommutative
$k$-group $G$ for which the functor of classes of pointed $G$-torsors
is not even a Zariski sheaf.
\end{proposition}

\begin{proof}
Let $k$ be algebraically closed of characteristic $3$.  Choose an
elliptic curve $C/k$ and a nontrivial line bundle $L$ of order $2$.
Such data are explicit: on $y^2=x^3-x$, the class
$L=\mathcal O_C((0,0)-(\infty))$ is nontrivial of order $2$, since
$\operatorname{div}(x)=2(0,0)-2(\infty)$.
Let $X$ be the split square-zero thickening
\[
\mathcal O_X=\mathcal O_C\oplus L,\qquad L^2=0
\text{ for this algebra multiplication}.
\]
As $L$ is nontrivial of degree zero,
$H^0(C,L)=0$ and hence $H^0(X,\mathcal O_X)=k$, universally after
base change over $k$.  The scheme is proper and flat over the field.
Choose the marked point through any $c\in C(k)$.

Let $Q_0\to C$ be the connected \'etale $\mu_2$-torsor associated
to the chosen isomorphism $L^{\otimes2}\simeq\mathcal O_C$, and let
$Q=Q_0\times_C X$.  Its global algebra is
$k[\epsilon]/(\epsilon^2)$, and the deck involution sends
$\epsilon$ to $-\epsilon$: the pullback of $L$ becomes trivial
with the nontrivial character.  Put
\[
G=\alpha_3\rtimes\mu_2,
\qquad (a,e)(b,f)=(a+eb,ef),
\]
and form the pointed $G$-torsor $P=Q\mathbin{\times^{\mu_2}}G$.

We calculate its pointed automorphism sheaf over arbitrary test
schemes.  For an affine test $T=\operatorname{Spec}R$, the
$\alpha_3$ part of a section of the adjoint group of $P_T$ is a
sign-equivariant function on $Q_T$.  Since
$H^0(Q_T,\mathcal O)=R\oplus R\epsilon$, these functions are
exactly $b\epsilon$, $b\in R$; they automatically have cube zero.
The $\mu_2$ part is a section of $\mu_2(T)$, and pointing forces it
to be $1$.  Every $b\epsilon$ vanishes at the marked point.  The
group law is addition of $b$.  Thus, functorially on $T$,
\[
\operatorname{Aut}(P_T,\text{pointing})=\mathbb G_a(T).
\]

Take a smooth proper curve $T/k$ of positive genus and a nonzero
class $\eta\in H^1(T,\mathcal O_T)=H^1(T,\mathbb G_a)$.
On a Zariski cover trivializing its additive torsor, glue the copies
of $P_T$ by the pointed automorphisms corresponding to a cocycle for
$\eta$.  Effective descent for finite flat torsors gives a pointed
$G_T$-torsor $P_\eta$ on $X_T$.  It is locally isomorphic to $P_T$.
If it were globally isomorphic as a pointed torsor, the cocycle
would be a coboundary in the full pointed automorphism sheaf
$\mathbb G_a$, contrary to $\eta\ne0$.  Hence two distinct elements
of the class functor become equal on a Zariski cover.  A representable
functor is a sheaf, so no scheme or algebraic space represents this
class functor.
\end{proof}

